\section{Reflection: approximated and traced}
\label{sec:rtadv}
\textit{Builds on: ray tracing (\ref{sec:rt}), illumination (\ref{sec:illum}).}

The raster renderer approximates mirror reflection with screen-space reflections (SSR): for a reflective pixel it
marches the reflected ray through the depth buffer and copies the colour where the ray passes behind a stored depth.
It can only reflect what is on screen. The hybrid mode replaces that with real reflection rays against a BVH of the
scene; Whitted, path and light tracing render the whole image with rays. The reflection weight is a Fresnel term.
\begin{equation}
F(\theta)=F_0+(1-F_0)(1-\cos\theta)^5\quad\text{(Schlick)},\qquad
L_{\text{out}}=(1-F)\,L_{\text{local}}+F\,L(\vect r)
\label{eq:fresnel}
\end{equation}
\begin{figure}[H]
\centering
\begin{subfigure}{0.49\linewidth}\raw{lab/reflections_compare}{\linewidth}\caption{Reflections off / SSR on (wet ground, metal).}\end{subfigure}\hfill
\begin{subfigure}{0.49\linewidth}\raw{rt/raytrace_ssr}{\linewidth}\caption{SSR vs ray-traced reflections.}\end{subfigure}
\caption{Approximate versus traced reflection. Notice in (b) where SSR fails: objects outside the screen or hidden
behind others are missing from the reflection, and the reflection fades towards the screen edge; the traced
reflection has them.}
\label{fig:refl}
\end{figure}

\begin{table}[H]
\centering\small
\rtperftable
\caption{Cost of each render mode at \rtperfsize{} on the RTX 4070 (\code{--capture raytrace\_perf}). Hybrid traces at full size; Whitted,
path and light tracing render at a reduced internal size and upscale.}
\label{tab:rtperf}
\end{table}

\begin{table}[H]
\centering\small
\begin{tabularx}{\linewidth}{@{}lXXl@{}}
\toprule
Method & Visual result & Cost & When to use\\\midrule
SSR & correct for visible geometry; holes for off-screen or hidden objects & one ray march per pixel in the depth buffer & real-time play (EPIC)\\
Hybrid RT reflections & complete reflections & BVH traversal per reflective pixel & stills, demonstration\\
Whitted & exact mirrors, refraction, hard shadows; no indirect light & recursive rays per pixel & teaching, specular scenes\\
Path / light tracing & soft shadows, indirect light, caustics (light) & many paths per pixel, noise, denoising & reference images\\
\bottomrule
\end{tabularx}
\caption{Differences between reflection methods.}
\end{table}
\where{SSR in src/rendering/ShadersPost.cpp; src/raytrace/RayTracer.cpp (Whitted, path, light tracing), Bvh.cpp; src/rendering/RayTracerGL.cpp}
\critique{The hybrid mode traces only the reflection. The screen-space passes (SSAO/GTAO, fog, SSR) work on the primary
image, so they cannot apply to what is seen in a traced reflection; Whitted, path and light tracing also run at reduced
internal resolution (Table~\ref{tab:rtperf}) and are for demonstration, not play.}
